% ===========================================================================
% Relatório — Projeto, Parte 2
% Segurança em Redes e Sistemas de Informação (SRSI)
%
% Compilar:   make            (pdflatex + bibtex + pdflatex x2)
% Dados:      make dados PI=<campanha> [DESKTOP=<campanha>]   (ver README.md)
% Usa o preâmbulo, as figuras e a bibliografia da Parte 1 (pasta acima).
% ===========================================================================
\documentclass[11pt, a4paper]{article}

\newcommand{\Mestrado}{Engenharia Informática}
\newcommand{\UnidadeCurricular}{Segurança em Redes e Sistemas de Informação}
\newcommand{\Autores}{%
  Rodrigo Gorjão, N.º 82972\\
  João Afonso, N.º 97990\\
  João Tendeiro, N.º 129087\\
  Jossivalter Leal, N.º 146851}
\newcommand{\Docente}{Daniel Maximino Caçador}
\newcommand{\Titulo}{Segurança em Camadas para Comunicações IoT Baseadas em MQTT}
\newcommand{\Subtitulo}{Implementação, Avaliação num Raspberry Pi e Validação}
\newcommand{\Entrega}{Projeto --- Parte 2: implementação e avaliação}
\newcommand{\TituloCurto}{Segurança em camadas para MQTT}
\newcommand{\Repositorio}{https://github.com/joaoaeafonso/mqtt_secure}

% ---------------------------------------------------------------------------
% Preâmbulo: pacotes, aspeto e macros do relatório.
% ---------------------------------------------------------------------------

\usepackage[T1]{fontenc}
\usepackage[utf8]{inputenc}
\usepackage[portuguese]{babel}
\usepackage{lmodern}
\usepackage{microtype}
\usepackage[a4paper, margin=2.5cm, headheight=14pt]{geometry}
\usepackage{setspace}
\onehalfspacing

% Matemática
\usepackage{amsmath, amssymb, mathtools, bm}

% Tabelas
\usepackage{array, booktabs, tabularx, longtable, multirow, makecell}
\usepackage{colortbl}
\usepackage{threeparttable}
\usepackage{pdflscape}

% Figuras
\usepackage{graphicx}
\usepackage{float}
\usepackage[font=small, labelfont=bf, labelsep=period]{caption}
\usepackage{subcaption}
\usepackage{tikz}
\usetikzlibrary{arrows.meta, positioning, shapes.geometric, shapes.multipart,
                fit, calc, backgrounds, decorations.pathreplacing}
\usepackage{pgfplots}
\pgfplotsset{compat=1.16}

% Diversos
\usepackage{xcolor}
\usepackage{enumitem}
\usepackage{pifont}
\usepackage{csquotes}
\usepackage{url}
\usepackage{fancyhdr}
\usepackage{titlesec}
\usepackage[hidelinks, unicode, pdfusetitle]{hyperref}

% ---------------------------------------------------------------------------
% Cores (inspiradas na identidade do Iscte; ajustáveis)
% ---------------------------------------------------------------------------
\definecolor{iscte}{RGB}{0, 51, 102}
\definecolor{acento}{RGB}{0, 112, 192}
\definecolor{cinza}{RGB}{110, 110, 110}
\definecolor{fundotab}{RGB}{235, 241, 248}
\definecolor{bom}{RGB}{0, 128, 0}
\definecolor{mau}{RGB}{190, 30, 45}
\definecolor{parcial}{RGB}{200, 130, 0}

% ---------------------------------------------------------------------------
% Títulos de secção e cabeçalhos
% ---------------------------------------------------------------------------
\titleformat{\section}{\Large\bfseries\color{iscte}}{\thesection}{0.8em}{}
\titleformat{\subsection}{\large\bfseries\color{iscte}}{\thesubsection}{0.8em}{}
\titleformat{\subsubsection}{\normalsize\bfseries\color{iscte}}{\thesubsubsection}{0.8em}{}

\pagestyle{fancy}
\fancyhf{}
\fancyhead[L]{\small\color{cinza}\TituloCurto}
\fancyhead[R]{\small\color{cinza}SRSI --- Projeto, Parte 1}
\fancyfoot[C]{\small\thepage}
\renewcommand{\headrulewidth}{0.4pt}

% Nomes em português de Portugal
\addto\captionsportuguese{%
  \renewcommand{\contentsname}{Índice}%
  \renewcommand{\listfigurename}{Lista de Figuras}%
  \renewcommand{\listtablename}{Lista de Tabelas}%
  \renewcommand{\refname}{Referências}%
  \renewcommand{\abstractname}{Resumo}%
  \renewcommand{\tablename}{Tabela}%
  \renewcommand{\figurename}{Figura}%
  \renewcommand{\appendixname}{Anexo}%
}

% ---------------------------------------------------------------------------
% Macros
% ---------------------------------------------------------------------------
% Marcas usadas nas tabelas comparativas
\newcommand{\Sim}{\textcolor{bom}{\ding{51}}}          % abordado
\newcommand{\Parc}{\textcolor{parcial}{\ding{108}}}    % abordado parcialmente
\newcommand{\Nr}{\textcolor{cinza}{--}}                % não reportado
\newcommand{\Nao}{\textcolor{mau}{\ding{55}}}          % ausente por desenho

% Identificadores dos estudos incluídos (E01..E27) com citação
\newcommand{\E}[1]{\textsf{E#1}}

% Colunas de tabela
\newcolumntype{L}[1]{>{\raggedright\arraybackslash}p{#1}}
\newcolumntype{C}[1]{>{\centering\arraybackslash}p{#1}}
\newcolumntype{Y}{>{\raggedright\arraybackslash}X}

% Caixas de definição/observação
\newcommand{\caixa}[2]{%
  \begin{center}\begin{tikzpicture}
    \node[draw=acento, fill=fundotab, rounded corners=2pt, inner sep=8pt,
          text width=0.9\linewidth, align=justify] {\textbf{#1.}\ #2};
  \end{tikzpicture}\end{center}}

\fancyhead[R]{\small\color{cinza}SRSI --- Projeto, Parte 2}
\hypersetup{pdftitle={\Titulo: \Subtitulo}, pdfsubject={SRSI -- Projeto, Parte 2}}
\setlength{\emergencystretch}{2em}

% ------------------------------------------------------------ macros -------
% Texto por escrever: visível no PDF, fácil de encontrar com grep.
\newcommand{\PorFazer}[1]{\textcolor{mau}{\textbf{[Por fazer:} #1\textbf{]}}}
% Referência à Parte 1 (documento separado: não há \ref entre os dois).
\newcommand{\PUm}[1]{Parte~1, #1}
% Tabela ou figura gerada por tools/analysis/analyze.py (make dados).
% \gerado{pi5/summary} lê gerado/pi5/summary.tex, ou mostra uma caixa se faltar.
\newcommand{\gerado}[1]{%
  \IfFileExists{gerado/#1.tex}{\input{gerado/#1}}{%
    \begin{center}\fbox{\parbox{0.9\linewidth}{\centering\small\color{mau}
      Em falta: \texttt{gerado/#1.tex}\\(\texttt{make dados} depois da campanha)}}\end{center}}}
% O mesmo para a plataforma de referência: as etiquetas ganham o prefixo "ref-",
% para não colidirem com as do Raspberry Pi.
\newcommand{\geradoRef}[1]{{%
  \let\etiquetaOriginal\label
  \renewcommand{\label}[1]{\etiquetaOriginal{ref-##1}}%
  \gerado{#1}}}
% Figura gerada, com legenda e etiqueta.
\newcommand{\figuraGerada}[3]{%
  \begin{figure}[ht]\centering\gerado{#1}\caption{#2}\label{#3}\end{figure}}

\makeatletter
\setlength{\@fptop}{0pt}
\setlength{\@fpbot}{0pt plus 1fil}
\makeatother

\begin{document}

\pagenumbering{Alph}
\begin{titlepage}
  \thispagestyle{empty}
  \centering
  \IfFileExists{../figuras/iscte-logo.png}%
    {\includegraphics[width=0.62\textwidth]{../figuras/iscte-logo.png}}{}
  \par\vspace{2.2cm}
  {\Large Mestrado em \Mestrado\par}
  \vspace{0.6cm}
  {\LARGE\bfseries\color{iscte} \UnidadeCurricular\par}
  \vspace{1.6cm}
  {\Large\bfseries \Titulo\par}
  \vspace{0.4cm}
  {\large \Subtitulo\par}
  \vspace{0.6cm}
  {\normalsize\itshape \Entrega\par}
  \vspace{1.6cm}
  \begin{flushleft}
    \large
    \textbf{Autores:}\\[0.2cm]
    \Autores\\[0.8cm]
    \textbf{Docente:} \Docente\\[0.8cm]
    \textbf{Repositório:} \url{\Repositorio}
  \end{flushleft}
  \vfill
\end{titlepage}

\pagenumbering{roman}
\section*{Resumo}
\addcontentsline{toc}{section}{Resumo}

A primeira parte deste projeto propôs uma arquitetura de segurança em camadas para
comunicações MQTT e um plano para a avaliar. Esta segunda parte apresenta o
protótipo que a implementa, as medições feitas num Raspberry~Pi~5 e a validação do
protótipo perante os ataques previstos. Cada camada pode ser ligada e desligada de
forma independente, o que permitiu medir o custo e o benefício de cada uma em doze
configurações (K0 a K11).

O transporte (mTLS) e o envelope AES-GCM custam cada um 7 a 8~$\mu$s de
processador por mensagem no \emph{gateway}, e as restantes verificações, juntas,
cerca de 2~$\mu$s; as assinaturas Ed25519 custam mais 151~$\mu$s e limitam o
débito a 5\,000~mensagens/s. O transporte e o envelope impedem, cada um, quase
todos os ataques de um estranho, mas os de um dispositivo comprometido só são
impedidos pela combinação de todas as camadas, que K8, K9 e K11 conseguem. A
configuração recomendada é K8 para a telemetria (30\,000~mensagens/s, p99 de
0,26~ms) e K9 para as ordens a atuadores, como a primeira parte previa. Das cinco
hipóteses, três confirmam-se; refuta-se que o custo do TLS esteja no
estabelecimento da ligação, e não se confirma que a limitação de tráfego poupe
processador quando a criptografia é AES em \emph{hardware}.

\medskip
\noindent\textbf{Palavras-chave:} MQTT, Internet das Coisas, segurança em camadas,
proteção fim-a-fim, avaliação de desempenho, Raspberry Pi.

\clearpage
\tableofcontents
\clearpage

\pagenumbering{arabic}
\section{Introdução}
\label{sec:introducao}

A primeira parte do projeto reviu a literatura sobre segurança em comunicações
MQTT, propôs uma arquitetura em camadas (\PUm{Secção~6}) e definiu o plano para a
avaliar (\PUm{Secção~7}). O objetivo desta segunda parte é o que esse plano
enuncia: medir, num Raspberry~Pi, quanto custa cada camada em desempenho e quanto
protege, isoladamente e em combinação, para identificar a configuração com a maior
proteção e o menor impacto no desempenho.

O documento responde às cinco perguntas experimentais da Parte~1:

\begin{enumerate}[label=\textbf{QP\arabic*.}, leftmargin=*]
  \item Qual é o custo de cada camada num Raspberry Pi (processador, latência,
    débito, memória, octetos)? (Secção~\ref{sec:custo})
  \item Que ataques impede cada camada, e quais só a combinação de camadas impede?
    (Secção~\ref{sec:protecao})
  \item Como varia o custo com o tamanho da mensagem e a carga, e quando satura o
    Raspberry Pi? (Secção~\ref{sec:custo})
  \item Quanto pesa a aceleração criptográfica por \emph{hardware}?
    (Secção~\ref{sec:custo})
  \item Que configurações são ótimas de Pareto, e qual recomendar para cada caso de
    uso? (Secção~\ref{sec:criterio})
\end{enumerate}

\paragraph{Contributos desta parte.}
\begin{itemize}[leftmargin=*]
  \item Um protótipo em C++23 da arquitetura proposta, com cada camada ativável de
    forma independente, e os seus testes automatizados (Secção~\ref{sec:prototipo}).
  \item Um laboratório de carga que executa de forma automática as doze
    configurações, os ataques e os respetivos controlos negativos contra um
    \emph{broker} real, e que guarda os dados brutos e o registo do ambiente de cada
    execução (Secção~\ref{sec:metodo}).
  \item As medições no Raspberry~Pi~5 e num computador de secretária, a aplicação
    do critério de escolha e o veredicto das hipóteses H1 a H5
    (Secções~\ref{sec:custo} a \ref{sec:hipoteses}).
  \item A validação do protótipo e a calibração dos limites de validade temporal
    que a Parte~1 deixou em aberto (Secção~\ref{sec:validacao}).
\end{itemize}

O plano da Parte~1 foi seguido; o que foi concretizado de forma diferente, ou
decidido só nesta parte, está declarado na Secção~\ref{sec:desvios}.

\section{Protótipo}
\label{sec:prototipo}

\subsection{Visão geral}

O protótipo tem três partes. A \textbf{biblioteca} (C++23) implementa as camadas
como elementos de uma cadeia de processamento: cada mensagem que chega passa pela
limitação de tráfego, pela camada de segurança (o envelope) e, opcionalmente, pela
persistência; cada camada pode ser ligada e desligada em execução, sem reconstruir
a cadeia. O \textbf{\emph{gateway} de ensaio} liga essa cadeia a um \emph{broker}
Mosquitto real através de um cliente Paho e reporta, por janela de medição, as
mensagens entregues e rejeitadas (por motivo), a latência e o tempo de processador.
O \textbf{laboratório de carga} (Python) lança o \emph{broker}, o \emph{gateway},
os sensores simulados e o atacante, aplica cada configuração e mede.

A Figura~\ref{fig:arquitetura} mostra onde corre cada parte na campanha.

\begin{figure}[ht]
  \centering
  \begin{tikzpicture}[
      proc/.style={draw=cinza, fill=white, rounded corners=2pt, font=\footnotesize,
                   align=center, minimum height=0.85cm, inner sep=2pt},
      camada/.style={draw=iscte, fill=fundotab, rounded corners=2pt, font=\footnotesize,
                     align=center, text width=4.1cm, minimum height=0.62cm, inner sep=2pt},
      maquina/.style={draw=iscte, thick, rounded corners=4pt, inner sep=6pt},
      rotulo/.style={font=\small\bfseries, text=iscte, anchor=south west},
      seta/.style={-{Stealth[length=2mm]}, thick, iscte},
      dupla/.style={{Stealth[length=2mm]}-{Stealth[length=2mm]}, thick, iscte},
      controlo/.style={-{Stealth[length=1.6mm]}, densely dashed, cinza},
      nota/.style={font=\scriptsize, text=cinza, align=center}]
    % --- computador de secretária
    \node[proc, text width=3.1cm] (s1) at (1.75, 0.9)
      {Agente de sensores 1\\\scriptsize 25 sensores e atuadores};
    \node[proc, text width=3.1cm] (s2) at (1.75, -0.4)
      {Agente de sensores 2\\\scriptsize 25 sensores e atuadores};
    \node[nota, below=0.08cm of s2] (env) {envelope em Python};
    \node[proc, text width=3.1cm, draw=mau] (atk) at (1.75, -2.3)
      {Agente atacante\\\scriptsize estranho ou dispositivo\\comprometido};
    \node[maquina, fit=(s1)(s2)(env)(atk)] (pc) {};
    \node[rotulo] at (pc.north west) {Computador de secretária};
    % --- Raspberry Pi 5
    \node[proc, text width=2.1cm, fill=iscte!10] (brk) at (6.25, 0.9)
      {\emph{Broker}\\Mosquitto 2.0};
    \node[nota, below=0.08cm of brk, text width=2.4cm] (brkn)
      {TLS/mTLS, autenticação, ACL, limites};
    \node[proc, text width=2.0cm, draw=mau, dashed] (rel) at (8.95, 0.9)
      {Intermediário\\\scriptsize só em ataques};
    \node[font=\scriptsize\bfseries, text=cinza] (gwt) at (12.85, 1.75)
      {\emph{Gateway} de ensaio (C++23)};
    \node[camada, fill=white, draw=cinza] (paho) at (12.85, 0.9) {Cliente MQTT (Paho)};
    \node[camada, below=0.2cm of paho] (lim) {Limitação de tráfego};
    \node[camada, below=0.2cm of lim] (sec) {\textbf{Segurança da mensagem}\\
      \scriptsize envelope, validade, janela,\\ políticas de modo e de chave};
    \node[camada, below=0.2cm of sec] (per) {Persistência (opcional)};
    \node[camada, fill=white, draw=cinza, below=0.2cm of per] (med)
      {Aplicação de ensaio\\\scriptsize entregas, rejeições,\\ latência, processador};
    \node[maquina, draw=cinza, thin, inner sep=4pt, fit=(gwt)(paho)(med)] (gw) {};
    \node[proc, text width=4.6cm] (ctl) at (7.6, -2.3)
      {Controlador do laboratório\\\scriptsize Python; configura, mede e guarda};
    \node[maquina, fit=(brk)(brkn)(rel)(gw)(ctl)] (pi) {};
    \node[rotulo] at (pi.north west) {Raspberry~Pi~5};
    % --- fluxos MQTT (medidos)
    \draw[dupla] (s1.east) -- ([yshift=0.2cm]brk.west);
    \draw[dupla] (s2.east) -- ([yshift=-0.05cm]brk.west);
    \draw[seta, mau] (atk.east) -- ([yshift=-0.3cm]brk.west);
    \draw[dupla] (brk) -- (rel);
    \draw[dupla] (rel) -- (paho);
    \foreach \a/\b in {paho/lim, lim/sec, sec/per, per/med}
      \draw[seta] (\a) -- (\b);
    % --- controlo do laboratório (fora do caminho medido)
    \draw[controlo] (ctl.west) -- (ctl.west -| pc.east);
    \draw[controlo] (ctl.north) -- (ctl.north |- brkn.south);
    \draw[controlo] (ctl.east) -- (ctl.east -| gw.west);
  \end{tikzpicture}
  \caption{Arquitetura implementada, tal como corre na campanha. A cadeia de camadas
    da Parte~1 (\PUm{Secção~6.1}) está no \emph{gateway} de ensaio. Setas
    contínuas: MQTT (telemetria dos sensores para o \emph{gateway}, ordens no
    sentido inverso; a vermelho, o atacante). O intermediário só existe nos ensaios
    de ataques na rede. A tracejado: o canal de controlo do laboratório (TLS), que
    lança e configura cada parte e fica fora do caminho medido.}
  \label{fig:arquitetura}
\end{figure}

\subsection{O que foi implementado}

A Tabela~\ref{tab:implementado} relaciona cada elemento da proposta com a sua
concretização.

\begin{table}[ht]
  \centering
  \caption{Elementos da arquitetura proposta e a sua implementação.}
  \label{tab:implementado}
  \footnotesize
  \begin{tabularx}{\textwidth}{@{}L{4.6cm}YL{4.2cm}@{}}
    \toprule
    \textbf{Elemento (Parte~1)} & \textbf{Implementação} & \textbf{Onde} \\
    \midrule
    Envelope de 25 octetos; tópico e modo nos dados associados; modos de cifra,
      autenticação e assinatura (§6.4) & AES-256-GCM, GMAC e Ed25519 sobre OpenSSL;
      a mesma implementação em Python para os sensores, verificada pelos mesmos
      vetores de teste & \texttt{modules/security}, \texttt{msec\_envelope.py} \\
    ChaCha20-Poly1305 como alternativa (§6.4) & O algoritmo é um atributo da chave:
      o formato do envelope não muda e um recetor serve os dois & \texttt{chacha20\_poly1305\_}\allowbreak\texttt{cipher.cpp} \\
    Ordem de verificação (§6.5) & Formato, modo aceite, validade, verificação
      criptográfica, chave--tópico, janela & \texttt{security\_layer.cpp} \\
    Validade temporal (§6.6) & $\Delta_a$ e $\Delta_s$ globais e por fluxo (filtro
      de tópico) & \texttt{FreshnessRule} \\
    Janela anti-repetição, $w = 64$ (§6.6) & Por chave, em memória, com estado
      limitado & \texttt{replay\_guard.cpp} \\
    Restrição chave--tópico (§6.7) & Política com padrões por dono da chave,
      aplicada depois da autenticação & \texttt{key\_policy.cpp} \\
    Chaves por época, rotação e revogação (§6.8) & Rotação com período de graça;
      revogação de chaves, palavras-passe e certificados (CRL) & \texttt{pki.py}, \texttt{lab.py} \\
    Limitação de tráfego antes da criptografia (§6.9) & \emph{Token bucket} por
      tópico & \texttt{modules/limits} \\
    Camada de transporte (§6.3) & TLS~1.3, autenticação mútua, palavras-passe,
      ACL por utilizador, identificador de cliente ligado ao utilizador, limites de
      pacote e de ligações & \texttt{broker.py} (Mosquitto) \\
    \bottomrule
  \end{tabularx}
\end{table}

\subsection{Decisões tomadas nesta parte}

\paragraph{Reinício do recetor.} A Parte~1 registou como limitação que a janela
anti-repetição se perde quando o recetor reinicia. A solução adotada não persiste a
janela a cada mensagem: o recetor guarda de forma durável apenas a época de chaves
mais recente que aceitou e, depois de um reinício, só aceita épocas posteriores. A
gravação faz-se uma vez por rotação, antes de aceitar a primeira mensagem da nova
época, e uma falha de leitura ou de escrita rejeita a mensagem.

\paragraph{Dependências entre mecanismos.} A validade temporal, a anti-repetição
e a restrição chave--tópico verificam campos do envelope. Uma configuração que as
ligue com o envelope desligado é recusada, no arranque e em execução, em vez de as
ignorar em silêncio.

\paragraph{Fora do âmbito.} Não foram implementados um modo HMAC próprio (o GMAC
cobre a autenticação sem cifra), \emph{nonces} por contador nem o modo de cifra
seguida de assinatura, que a Parte~1 apresenta como possível extensão.

\subsection{Testes automatizados}

O protótipo tem 93 testes em C++
(unitários, de integração e de aceitação), executados com AddressSanitizer e
UndefinedBehaviorSanitizer, e 73 testes em Python (61 do laboratório, 12 da
análise). Os testes de aceitação têm um
teste por regra da ordem de verificação e foram verificados por mutação: cada um
falha quando a regra correspondente é removida. Três alvos de \emph{fuzzing}
exercitam a análise do envelope, a validação de tópicos e todo o caminho de
receção com octetos arbitrários. A integração contínua compila com
\texttt{clang-tidy} (avisos tratados como erros) e corre os cenários de ataque
contra um Mosquitto real.

\section{Método experimental}
\label{sec:metodo}

\subsection{Plataforma}

\PorFazer{Preencher a partir de \texttt{manifest.json} de cada campanha: modelo e
memória do Raspberry~Pi~5, sistema operativo e \emph{kernel}, versões do OpenSSL,
do Mosquitto e do compilador, \emph{commit} do código; o mesmo para o computador de
secretária e para as máquinas dos sensores e do atacante; rede (cabo ou Wi-Fi);
dissipador; regulador de frequência; se o NTP esteve ativo.}

Tal como planeado, a avaliação tem dois níveis. No \textbf{nível~1}, os
\emph{microbenchmarks} medem o tempo de processador da cadeia de camadas por
mensagem, sem rede, para 16~B, 256~B, 4~KiB e 64~KiB. No \textbf{nível~2}, o
\emph{broker} e o \emph{gateway} correm no Raspberry~Pi e os sensores e o atacante
noutras máquinas da rede local.

\subsection{Configurações}

Foram medidas as doze configurações da \PUm{Tabela das configurações} (K0 a K11),
sem alterações. O nível~1 tem um perfil para K0, K2 a K7, K9 e K11; K1, K8 e K10
só diferem das vizinhas no transporte e nos limites, que só existem no nível~2.

\subsection{Métricas}

A Tabela~\ref{tab:medidas} indica como cada métrica do plano foi medida.

\begin{table}[ht]
  \centering
  \caption{Como cada métrica foi medida.}
  \label{tab:medidas}
  \footnotesize
  \begin{tabularx}{\textwidth}{@{}L{3.8cm}Y@{}}
    \toprule
    \textbf{Métrica} & \textbf{Medição} \\
    \midrule
    Processador por mensagem & Tempo de processador do processo do \emph{gateway}
      na janela, dividido pelas mensagens recebidas (nível~2); tempo de processador
      por mensagem da cadeia completa (nível~1) \\
    Latência p50 e p99 & Do instante de envio, escrito pelo sensor na mensagem, à
      entrega depois da cadeia de camadas. O relógio de cada máquina de sensores é
      corrigido pelo desvio estimado em relação ao do \emph{gateway} \\
    Débito máximo sustentado & Maior taxa entregue com perda inferior a 1\,\% e
      p99 dentro de $L_{\max}$, numa subida de carga por configuração \\
    Octetos por mensagem & Octetos recebidos na ligação TCP do \emph{gateway} por
      mensagem, separados em conteúdo, envelope, MQTT e restante (TLS e controlo) \\
    Memória & Pico de memória residente do \emph{gateway} e do \emph{broker} na
      janela \\
    Estabelecimento de ligação & Tempo e tempo de processador de religações do
      \emph{gateway} (TLS, CONNECT e subscrições), fora da janela de tráfego \\
    Ataques impedidos, $\pi_a(c)$ & Cada ataque contra cada configuração, com
      tráfego legítimo ativo; maioria sobre as repetições \\
    Fuga de informação & Fração das mensagens legíveis por um subscritor de
      \texttt{\#} e por um atacante na ligação do \emph{gateway} \\
    \bottomrule
  \end{tabularx}
\end{table}

A energia, opcional no plano, não foi medida. \PorFazer{Confirmar, ou descrever o
medidor se vier a ser usado.}

\subsection{Ataques}

O atacante do laboratório executa os ataques da \PUm{Tabela dos ataques}. Dois
tipos de ensaio usam-nos de forma diferente:

\begin{itemize}[leftmargin=*]
  \item A \textbf{matriz de ataques} lança cada ataque contra cada configuração
    completa e dá $\pi_a(c)$, de que resulta $S(c)$. Um ataque travado no
    \emph{broker} (por exemplo, um estranho sem credenciais perante a autenticação
    mútua) conta como impedido por essa configuração.
  \item Os \textbf{cenários} isolam uma defesa de cada vez sobre uma base aberta e
    emparelham cada ataque com um \emph{controlo negativo}, a mesma configuração sem
    essa defesa, em que o ataque tem de ter sucesso. Só assim se atribui uma
    rejeição à camada certa.
\end{itemize}

Os ataques internos partem de um dispositivo comprometido: o atacante recebe as
chaves, a palavra-passe e o certificado de um sensor, e o sensor verdadeiro fica
desligado enquanto o atacante o controla. Os ataques na rede usam um intermediário
na ligação entre o \emph{broker} e o \emph{gateway}, que conta o que consegue ler
e, em MQTT sem TLS, altera, redireciona, atrasa ou suprime mensagens, fazendo de
\emph{broker} malicioso.

\subsection{Procedimento}

Cada configuração foi repetida \PorFazer{10} vezes, em ordem aleatória com semente
registada, com 3~s de estabilização e 30~s de medição por execução. Os intervalos
de confiança são de 95\,\%, por \emph{bootstrap} de percentis sobre as execuções.
O custo marginal de uma camada é a diferença de medianas entre as duas
configurações que só diferem nela, e a diferença das diferenças mede a interação
entre a validade temporal e a anti-repetição e entre o TLS e o envelope. A matriz
de ataques correu separada das medições de custo, porque o intermediário na
ligação do \emph{gateway} acrescenta latência.

Cada campanha guarda os dados brutos de todas as execuções, a ordem dos ensaios, a
semente, a configuração completa e um registo do ambiente (código, plataforma,
versões, temperatura e frequência do processador por execução). Todas as tabelas e
figuras deste relatório são geradas a partir desses dados por um único comando
(Anexo~\ref{sec:reproducao}).

\subsection{Parâmetros fixados antes das medições}
\label{sec:parametros}

Os parâmetros do critério de escolha foram fixados antes da campanha: o ficheiro
que os contém é copiado para a pasta da campanha antes da primeira medição e o seu
resumo criptográfico fica no registo do ambiente.

\begin{itemize}[leftmargin=*]
  \item \textbf{Pesos dos ataques.} Cada linha da tabela de ataques da Parte~1
    pesa~1. Duas linhas são medidas por dois ataques cada (a escuta, na rede e no
    \emph{broker}; a injeção e a adulteração), que pesam 0,5 cada.
  \item \textbf{Pesos do custo.} Processador por mensagem 0,4; latência p99 0,4;
    débito máximo, octetos por mensagem, memória do \emph{gateway} e
    estabelecimento de ligação 0,05 cada.
  \item \textbf{Exequibilidade.} $R_{\text{alvo}} = $ \PorFazer{valor}~mensagens/s
    e $L_{\max} = $ \PorFazer{valor}~ms, escolhidos depois de um piloto no
    Raspberry~Pi (\PorFazer{justificar com o cenário de aplicação}) e antes da
    campanha.
\end{itemize}

\subsection{Desvios ao plano}
\label{sec:desvios}

\begin{itemize}[leftmargin=*]
  \item \textbf{Aceleração AES desligada.} O plano previa
    \texttt{OPENSSL\_armcap=0}. Essa variável desliga todas as extensões ARM que o
    OpenSSL usa, incluindo o NEON, de que o ChaCha20-Poly1305 depende. Um
    dispositivo sem as extensões criptográficas tem, em geral, NEON; por isso
    mediu-se também com \texttt{OPENSSL\_armcap=1} (só NEON), e é essa a
    comparação usada para a QP4. Os dois resultados são apresentados.
  \item \textbf{Estabelecimento de ligação.} O tempo de ligação medido varia em
    degraus de cerca de 100~ms, por esperas internas da biblioteca cliente, que
    escondem os poucos milissegundos do TLS. No custo agregado entra o tempo de
    processador por ligação; o tempo decorrido é reportado à parte.
  \item \textbf{Escuta no \emph{broker}.} Na matriz de ataques, este ataque é um
    subscritor de \texttt{\#} sem credenciais. Um \emph{broker} curioso não é
    medido por $S(c)$ (Secção~\ref{sec:ameacas}); é o requisito de
    confidencialidade do caso de uso que o cobre.
  \item \textbf{Dispositivo comprometido.} O sensor verdadeiro fica desligado
    durante o ataque. Com os dois ligados, disputam a sessão MQTT e o ataque fica
    reduzido a uma fração do tempo (Secção~\ref{sec:observacoes}).
  \item \PorFazer{Outros desvios que surjam na campanha (repetições a mais,
    execuções repetidas, energia).}
\end{itemize}

\section{Resultados: custo}
\label{sec:custo}

\subsection{Custo por configuração (QP1)}

As Tabelas~\ref{tab:resumo} e~\ref{tab:resumo-b} dão, para cada configuração, a
mediana e o intervalo de confiança de cada métrica no Raspberry~Pi~5.

\gerado{pi5/summary}

\figuraGerada{pi5/fig-cpu}{Tempo de processador do \emph{gateway} por mensagem,
  por configuração (mediana e IC 95\,\%).}{fig:cpu}
\figuraGerada{pi5/fig-p99}{Latência p99 por configuração (mediana e IC
  95\,\%).}{fig:p99}

A 1\,000~mensagens/s, as configurações formam quatro grupos
(Figuras~\ref{fig:cpu} e~\ref{fig:p99}). K0, sem proteção, gasta
24,0~$\mu$s de processador do \emph{gateway} por mensagem, com latência p50 de
122~$\mu$s e p99 de 173~$\mu$s. O transporte sozinho (K1) e o envelope sozinho
(K10) ficam a meio: 31,9 e 35,2~$\mu$s, p99 de 220 e 224~$\mu$s. As
configurações com transporte e envelope AES-GCM (K2 a K8 e K11) formam um só
grupo, entre 39,1 e 42,6~$\mu$s e com p99 entre 257 e 269~$\mu$s; dentro dele,
as diferenças são de 1 a 3~$\mu$s e quase todas cabem no ruído. K9, com
assinaturas Ed25519, está isolada: 192,6~$\mu$s por mensagem, oito vezes K0, e p99
de 726~$\mu$s. Mesmo em K9, a latência continua abaixo de 1~ms, cinquenta vezes
menos do que $L_{\max}$; à carga nominal, o custo de todas as configurações é
pequeno em termos absolutos, e o que as separa é o débito máximo
(Secção~\ref{sec:saturacao}).

\subsection{Custo marginal de cada camada}

\gerado{pi5/marginal}

A Tabela~\ref{tab:marginal} dá o custo de cada camada isolada:

\begin{itemize}[leftmargin=*]
  \item \textbf{Transporte} (K0$\to$K1): +7,9~$\mu$s por mensagem e +47~$\mu$s
    no p99, mais 22~octetos por mensagem na ligação do \emph{gateway} e 4,8~MB de
    memória no \emph{broker}.
  \item \textbf{Envelope} (K1$\to$K2): +7,3~$\mu$s e +41~$\mu$s no p99, quase
    tanto como o TLS, e 53~octetos por mensagem. Passar de GMAC a AES-256-GCM
    (K2$\to$K3) não custa nada mensurável ($-0{,}1$~$\mu$s, IC $[-0{,}5; 0{,}5]$):
    com AES em \emph{hardware}, cifrar 16~octetos é desprezável face ao resto do
    envelope (análise, procura da chave, verificações).
  \item \textbf{Validade temporal} (K3$\to$K4) e \textbf{anti-repetição}
    (K3$\to$K5): diferenças de 0,6 e 0,2~$\mu$s, com intervalos que incluem zero.
  \item \textbf{Restrição chave--tópico} (K6$\to$K7): +1,3~$\mu$s
    (IC $[0{,}7; 1{,}8]$), a única verificação barata que se distingue do ruído.
  \item \textbf{Limitação de tráfego e limites do \emph{broker}} (K7$\to$K8):
    +0,4~$\mu$s, dentro do ruído.
  \item \textbf{Assinatura em vez de AES-GCM} (K8$\to$K9): +151,5~$\mu$s por
    mensagem, +464~$\mu$s no p99 e +36~octetos (89 em vez de 53).
  \item \textbf{ChaCha20-Poly1305 em vez de AES-256-GCM} (K8$\to$K11): +1,5~$\mu$s,
    num processador com AES em \emph{hardware} (Secção~\ref{sec:aceleracao}).
\end{itemize}

\paragraph{Interações.} A validade e a anti-repetição custam o mesmo juntas e
separadas: a diferença das diferenças é de $-0{,}5$~$\mu$s, IC $[-1{,}3; 0{,}4]$.
O TLS custa menos com o envelope ativo do que sem ele (5,5 contra 7,9~$\mu$s, uma
interação de $-2{,}4$~$\mu$s, IC $[-2{,}9; -2{,}0]$); a causa não foi
investigada. Para a latência, nenhuma das interações se distingue do ruído.

\subsection{Tamanho da mensagem e saturação (QP3)}
\label{sec:saturacao}

\paragraph{Tamanho.} No nível~1 (Tabela~\ref{tab:microbench}), sem rede, o custo de
K0 quase não depende do tamanho (0,8~$\mu$s a 16~B, 7,1~$\mu$s a 64~KiB). Com AES
em \emph{hardware}, o envelope AES-GCM (K3) custa 5,6~$\mu$s a 16~B, 6,0 a 256~B,
13,0 a 4~KiB e 118~$\mu$s a 64~KiB: até alguns KiB, domina a parte fixa do
envelope, e só acima disso o custo passa a crescer com o tamanho. A assinatura
(K9) é dominada pela parte fixa em toda a gama (266~$\mu$s a 16~B, 301~$\mu$s a
4~KiB, 820~$\mu$s a 64~KiB). No nível~1 cada mensagem é selada e aberta, pelo que
estes valores incluem o emissor e o recetor; no nível~2 o \emph{gateway} só
verifica, e a diferença K8$\to$K9 é de 151~$\mu$s. O tamanho não foi variado no
nível~2.

\paragraph{Saturação.} A Tabela~\ref{tab:saturacao} dá o débito máximo sustentado
(perda inferior a 1\,\% e p99 até $L_{\max}$) e o processador no último patamar
sustentado. Em todas as configurações, o que limita primeiro é o processador do
\emph{gateway}: no patamar seguinte, o seu processo chega a 93--100\,\% de um núcleo
(a cadeia de camadas corre numa só linha de execução), enquanto o \emph{broker}
fica entre 15 e 80\,\%. K0 sustenta 60\,000~mensagens/s, o transporte desce para
40\,000, as configurações com envelope AES ficam entre 25\,000 e 35\,000 e K9 fica
em 5\,000~mensagens/s, metade de $R_{\text{alvo}}$ da telemetria. Os patamares são
discretos (passos de 5\,000~mensagens/s entre 400 e 800 por sensor): uma
diferença de um patamar entre K3 e K8 está dentro da resolução da subida. K6, K7
e K10 sustentaram 30\,000~mensagens/s numa das três subidas e falharam nas
outras duas, com p99 de 120 a 520~ms; a mediana fica no patamar anterior.

\begin{table}[ht]
  \centering
  \caption{Débito máximo sustentado no Raspberry~Pi~5 (mediana das três subidas)
    e processador do \emph{gateway} e do \emph{broker} no último patamar
    sustentado.}
  \label{tab:saturacao}
  \footnotesize
  \begin{tabular}{@{}lrrr@{}}
    \toprule
    \textbf{Config.} & \textbf{Débito (msg/s)} & \textbf{\emph{Gateway} (\%)} &
      \textbf{\emph{Broker} (\%)} \\
    \midrule
    K0  & 60\,000 & 95 & 56 \\
    K1  & 40\,000 & 93 & 72 \\
    K2  & 35\,000 & 96 & 61 \\
    K3, K4, K5 & 30\,000 & 84--87 & 55--57 \\
    K6, K7 & 25\,000 & 74--76 & 44--48 \\
    K8  & 30\,000 & 89 & 55 \\
    K9  & 5\,000  & 93 & 14 \\
    K10 & 25\,000 & 70 & 45 \\
    K11 & 30\,000 & 92 & 57 \\
    \bottomrule
  \end{tabular}
\end{table}

A partir de 400~mensagens/s por sensor (20\,000~mensagens/s), a latência sobe de
forma brusca em todas as configurações, incluindo K0: o p99 passa de 0,2--0,6~ms
para 20--35~ms e depois volta a descer com a carga, sem perda. Isto acontece com o
\emph{gateway} a 40\,\% (K0) ou a 63\,\% (K8) de processador, e K9 satura antes de
lá chegar. O efeito foi registado, mas a sua origem não foi determinada
(Secção~\ref{sec:ameacas}). Fica abaixo de $L_{\max}$ e acima de
$R_{\text{alvo}}$ (200~mensagens/s por sensor), pelo que não muda nenhum veredicto.

\subsection{Aceleração criptográfica por \emph{hardware} (QP4)}
\label{sec:aceleracao}

\gerado{pi5/microbench}

O Cortex-A76 do Raspberry~Pi~5 tem as extensões criptográficas ARMv8, e o
OpenSSL usa-as para o AES e o GHASH. Sem elas, só com NEON
(\texttt{OPENSSL\_armcap=1}), o envelope AES-GCM (K3) passa de 5,6 para 7,5~$\mu$s
a 16~B, de 6,0 para 12,4 a 256~B e de 13,0 para 88~$\mu$s a 4~KiB, sete vezes
mais. O ChaCha20-Poly1305 (K11) quase não depende das extensões (6,9 a 7,6~$\mu$s
a 16~B; 21,7 a 22,4~$\mu$s a 4~KiB), porque usa NEON; só sem NEON
(\texttt{OPENSSL\_armcap=0}) chega a 38~$\mu$s a 4~KiB.

Daqui resulta a regra de escolha da Parte~1, agora com números. Com AES em
\emph{hardware}, o AES-GCM é o mais barato em todos os tamanhos (13,0 contra
21,7~$\mu$s a 4~KiB), e no nível~2 K11 custa mais 1,5~$\mu$s por mensagem do que
K8. Sem AES em \emph{hardware}, o ChaCha20 compensa a partir de algumas centenas
de octetos: a 256~B já custa 8,7 contra 12,4~$\mu$s, e a 4~KiB custa quatro vezes
menos. A 16~B, a diferença é inferior a 0,1~$\mu$s e quem decide é o resto da
cadeia. A assinatura não muda com as extensões, porque o Ed25519 não as usa.

\subsection{Octetos e memória}

O envelope acrescenta 53~octetos por mensagem com AES-GCM ou ChaCha20-Poly1305 e
89 com Ed25519, exatamente o que o formato prevê; com o conteúdo de 16~octetos
usado na campanha, a mensagem na ligação do \emph{gateway} passa de 102 octetos
(K0) para 178 (K2 a K8, K11) e 214 (K9). O TLS acrescenta 22~octetos por
mensagem, a que se soma a negociação de cada ligação, que custa 4,2~ms de
processador no \emph{gateway} com autenticação mútua, contra 0,8~ms sem TLS. A
memória do \emph{gateway} é a mesma em todas as configurações (cerca de 39~MB, um
majorante, Secção~\ref{sec:ameacas}); a do \emph{broker} passa de 7,5~MB sem TLS
para cerca de 12,3~MB com TLS, e não muda com as restantes camadas.

\subsection{Plataforma de referência}

\PorFazer{O mesmo no computador de secretária (Anexo~\ref{sec:reproducao}): em que
é que a ordenação das configurações muda, se mudar.}

\section{Resultados: proteção}
\label{sec:protecao}

\subsection{Ataques impedidos por configuração (QP2)}

\PorFazer{Tabela $\pi_a(c)$: configurações em linhas, ataques em colunas, a partir
de \texttt{gerado/pi5/protection.csv}, com $S(c)$ na última coluna. Indicar o
número de repetições e os casos em que as repetições não concordaram.}

\PorFazer{Leitura: o que K1 já impede (tudo o que exige chegar ao \emph{broker});
o que só o envelope impede (K10 sem TLS, \emph{broker} ativo); o que só a
combinação impede (mensagem antiga nunca recebida: validade; repetição: janela);
o que só K8 e seguintes impedem (inundação, ocupação de ligações).}

\subsection{Cada ataque contra a sua defesa, com controlo negativo}

\PorFazer{Tabela dos cenários (\texttt{results/<campanha>/scenarios.md}): ataque,
defesa, previsão, resultado e evidência, com o controlo negativo ao lado. Esta
tabela confirma ou refuta, linha a linha, a tabela de ataques da Parte~1.}

\subsection{Observações}
\label{sec:observacoes}

Os ensaios mostraram efeitos que a tabela de ataques da Parte~1 não previa.
\PorFazer{Confirmar cada um com os dados do Raspberry~Pi e quantificar.}

\begin{itemize}[leftmargin=*]
  \item \textbf{Ocupação de ligações.} O limite de ligações protege os recursos do
    \emph{broker} e mantém os clientes já ligados, mas o atacante ocupa as vagas
    livres; se o \emph{broker} reiniciar durante o ataque, o atacante pode ficar
    com todas as vagas e nenhum sensor volta a ligar-se.
  \item \textbf{Chave roubada e janela da vítima.} As mensagens aceites com uma
    chave roubada avançam a janela anti-repetição dessa chave; o tráfego legítimo
    do dispositivo passa a ser rejeitado até à rotação de chaves.
  \item \textbf{Credencial clonada e sessão MQTT.} Se o atacante se liga com a
    identidade de um dispositivo que continua ligado, o \emph{broker} entrega-lhe
    a sessão e desliga o legítimo; os dois alternam enquanto o ataque dura, o que
    é uma negação de serviço da vítima.
  \item \textbf{Falsificação pelo recetor num \emph{broker} legítimo.} A ACL já
    impede um atuador de publicar nos tópicos de ordens dos outros; a diferença
    entre a chave partilhada e a assinatura só se observa sem ACL ou com um
    \emph{broker} ativo.
  \item \textbf{Atraso e supressão.} Um \emph{broker} ativo que atrase mensagens
    menos do que $\Delta_a$, ou que as suprima, não é detetado pelo recetor.
\end{itemize}

\section{Critério de escolha e recomendação}
\label{sec:criterio}

Aplica-se o critério da \PUm{Secção~7.5} com os parâmetros da
Secção~\ref{sec:parametros}: proteção $S(c)$, custo normalizado por K0,
exequibilidade, requisitos por caso de uso e fronteira de Pareto.

\gerado{pi5/criterion}

\figuraGerada{pi5/fig-pareto}{Proteção $S(c)$ contra custo$(c)$; a tracejado, a
  fronteira de Pareto para a telemetria sensível.}{fig:pareto}

\PorFazer{Para cada caso de uso (telemetria não sensível, telemetria sensível,
ordens a atuadores): configurações admissíveis, fronteira de Pareto, $c^*$ e custo
marginal $\Delta\operatorname{custo}/\Delta S$ ao longo da fronteira. Confirmar se
a recomendação antecipada na Parte~1 (K8, ou K11 sem AES em \emph{hardware}, para a
telemetria; K9 para as ordens) se mantém. Sensibilidade: a recomendação muda com
pesos diferentes?}

\section{Hipóteses}
\label{sec:hipoteses}

A evidência de cada hipótese está em \texttt{gerado/pi5/hypotheses.md}; o
veredicto é dos autores.

\begin{table}[ht]
  \centering
  \caption{Veredicto das hipóteses da Parte~1.}
  \label{tab:hipoteses}
  \footnotesize
  \begin{tabularx}{\textwidth}{@{}L{5.6cm}YL{2.2cm}@{}}
    \toprule
    \textbf{Hipótese} & \textbf{Evidência} & \textbf{Veredicto} \\
    \midrule
    H1. O envelope AES-256-GCM tem um custo por mensagem pequeno face à latência
      da rede & K1$\to$K3: +7,2~$\mu$s de processador (IC $[6{,}8; 7{,}8]$) e
      +25~$\mu$s no p50 (149~$\mu$s em K1); débito máximo de 40\,000 para
      30\,000~mensagens/s & Confirmada, com ressalva \\
    H2. Validade, anti-repetição e chave--tópico têm custo marginal quase nulo &
      K3$\to$K4 +0,6~$\mu$s e K3$\to$K5 +0,2~$\mu$s (IC incluem 0); K6$\to$K7
      +1,3~$\mu$s (3\,\%); p99 sem diferença mensurável & Confirmada \\
    H3. Ed25519 custa claramente mais do que AES-GCM por mensagem &
      K8$\to$K9: +151,5~$\mu$s (4,7 vezes), +464~$\mu$s no p99, débito máximo seis
      vezes menor (5\,000 contra 30\,000~mensagens/s) & Confirmada \\
    H4. O custo do TLS concentra-se no estabelecimento da ligação &
      Por ligação: 4,2 contra 0,8~ms de processador. Por mensagem: K0$\to$K1
      +7,9~$\mu$s e K10$\to$K7 +5,5~$\mu$s; débito máximo de 60\,000 para
      40\,000~mensagens/s & Refutada \\
    H5. Sob inundação, a limitação antes da criptografia reduz o processador do
      \emph{gateway} & Inundação a 300~msg/s: K7 5,2\,\%, K8 5,1\,\% de
      processador (40,5 e 39,7~$\mu$s por mensagem); K8 descarta 1\,559
      mensagens antes da criptografia & Não confirmada \\
    \bottomrule
  \end{tabularx}
\end{table}

\paragraph{H1 (confirmada, com ressalva).} Em valor absoluto, o envelope custa
pouco: 7,2~$\mu$s de processador e 25~$\mu$s de latência mediana, mil vezes menos
do que $L_{\max}$ e menos do que a latência de qualquer rede sem fios típica da
IoT. Mas, nesta rede local cablada, a latência é tão baixa que o envelope
representa 17\,\% do p50 de K1, e reduz o débito máximo em 25\,\%. A hipótese
vale para a latência; para o débito, o envelope não é desprezável.

\paragraph{H4 (refutada).} A negociação TLS com autenticação mútua custa
3,4~ms de processador a mais por ligação, mas o TLS custa também 7,9~$\mu$s por
mensagem, tanto como o envelope inteiro (7,2~$\mu$s), e é o que mais reduz o
débito máximo (de 60\,000 para 40\,000~mensagens/s). Ao fim de cerca de 430
mensagens, o custo por mensagem ultrapassa o da negociação; a 1\,000~mensagens/s,
isso acontece em menos de meio segundo, e um sensor a 10~Hz chega lá em 43~s. Em
ligações duradouras, que são as do MQTT, o custo do TLS está nas mensagens, não
no estabelecimento. O custo por mensagem é menor com o envelope ativo (5,5 contra
7,9~$\mu$s, Secção~\ref{sec:custo}), mas continua a dominar.

\paragraph{H5 (não confirmada).} A limitação funciona como previsto: descarta o
excesso antes de qualquer verificação criptográfica (1\,559 de 1\,814 mensagens
da inundação). Mas o processador do \emph{gateway} é o mesmo em K7 e K8 (5,2 e
5,1\,\%). Com AES em \emph{hardware}, verificar uma mensagem custa poucos
microssegundos: o que a limitação poupa, cerca de 260 verificações por segundo,
é da ordem de 0,2\,\% de um núcleo, abaixo do ruído. Com verificações mais caras,
como as assinaturas, a poupança poderia ser maior, mas a matriz não tem uma
configuração com assinaturas e sem limitação que o meça. A hipótese, tal como formulada para o
envelope AES-GCM, não se confirma. O que a limitação mostrou foi a proteção da
aplicação: sem ela (K7), as mensagens da inundação são todas entregues; com ela
(K8), o excesso é descartado (Secção~\ref{sec:protecao}).

\section{Validação do protótipo}
\label{sec:validacao}

\subsection{Conformidade com a especificação}

A validação tem três níveis. Os \textbf{testes automatizados}
(Secção~\ref{sec:prototipo}) verificam cada regra de aceitação, os vetores de
teste do envelope (partilhados pela implementação em C++ e pela em Python, que
assim não podem divergir) e os vetores do RFC~8439 para o ChaCha20-Poly1305. Os
\textbf{cenários} (Secção~\ref{sec:protecao}) verificam o sistema completo contra
um \emph{broker} real: os 39 cenários deram o resultado previsto no
Raspberry~Pi, na própria campanha (\texttt{results/pi5-20261113/scenarios.md}). No
mesmo \emph{commit} da campanha (\texttt{79c4878}), os 93 testes em C++ e os 61
testes em Python do laboratório passaram no Raspberry~Pi
(\texttt{results/validation/pi5-20261113/}); só os sete testes que dependem do
\texttt{tshark} não correram, por o Wireshark não estar instalado no
Raspberry~Pi.
A \textbf{matriz de ataques} confirma que cada configuração impede o que o desenho
prevê (Tabela~\ref{tab:protecao}), com as três repetições de acordo em todas as
células.

\subsection{Validação do instrumento de medição}

Os pilotos no Raspberry~Pi validaram também o próprio laboratório, e revelaram
dois defeitos que teriam falseado a campanha:

\begin{itemize}[leftmargin=*]
  \item \textbf{Deriva dos relógios.} No primeiro piloto, a latência mediana de K0
    saiu negativa ($-455~\mu$s) e variava de 300 a 1\,055~$\mu$s entre repetições
    da mesma configuração. O desvio entre os relógios crescia de forma regular
    (de 32,9 para 51,8~ms ao longo da subida de carga, cerca de 13~ppm), e a
    estimação usava a melhor amostra do último minuto, que podia estar 0,8~ms
    atrasada. Com a reta ajustada às amostras
    (Secção~\ref{sec:metodo}), a latência mediana de K0 ficou entre 132 e
    133~$\mu$s nas três repetições do terceiro piloto. Um teste reproduz a
    deriva observada e exige um erro inferior a 50~$\mu$s.
  \item \textbf{Contaminação entre patamares.} No segundo piloto, uma execução
    de K8 no primeiro patamar (500~mensagens/s) registou 77\,388 rejeições e p99
    de 1,3~s: seguia-se a um patamar de K11 em sobrecarga, cujas mensagens retidas
    chegaram depois da mudança de chave e de limite de débito. Se a perda tivesse
    passado de 1\,\%, a subida de carga de K8 teria parado logo nesse patamar e o
    seu débito máximo ficaria por medir. Com a espera
    pelo escoamento, nenhuma execução posterior a uma saturação teve rejeições no
    terceiro piloto.
\end{itemize}

\subsection{Calibração da validade temporal}

A Parte~1 fixou valores iniciais, $\Delta_a = 60$~s e $\Delta_s = 5$~s, a calibrar
nesta parte a partir da latência medida e do desvio de relógio observado. Cada
execução regista a cauda da latência e o desvio entre o relógio de cada máquina de
sensores e o do \emph{gateway}.

Na campanha, a maior latência de uma mensagem legítima foi de 43,2~ms (p99,9
máximo de 37,1~ms), contando as execuções de custo e os patamares da subida de
carga que cada configuração sustenta. Com o NTP desligado, o desvio entre os
relógios das duas máquinas cresceu até 289~ms ao longo das cinco horas; depois da
correção pelos agentes, a incerteza residual nunca passou de 0,21~ms. Com margem
de segurança~2 e arredondamento a 100~ms, os valores propostos são:

\begin{itemize}[leftmargin=*]
  \item \textbf{relógios sincronizados} (NTP ou PTP nos dispositivos):
    $\Delta_s = 100$~ms e $\Delta_a = 100$~ms;
  \item \textbf{relógios sem sincronização}, em que $\Delta_s$ tem de cobrir o
    desvio observado: $\Delta_s = 600$~ms e $\Delta_a = 700$~ms.
\end{itemize}

Os dois casos estão muito abaixo dos 60~s e 5~s iniciais da Parte~1: uma mensagem
antiga só é aceite se tiver menos de um segundo, em vez de um minuto. O segundo
caso depende de quanto tempo os relógios ficam sem acerto (aqui, cinco horas a
cerca de 16~ppm); num dispositivo sem relógio de tempo real que reinicia, o desvio
pode ser arbitrário, e é a sincronização que o resolve, não um $\Delta_s$ maior.
Para as ordens a atuadores, em que uma ordem atrasada é ela própria um risco, a
validade curta é a que interessa: com relógios sincronizados, uma ordem com mais
de 100~ms é recusada. Por fluxo (\texttt{FreshnessRule}), a telemetria pode
manter uma validade mais larga sem afetar as ordens.

\subsection{Reinício do recetor}

O recetor guarda em disco a última época de chaves que aceitou. Depois de um
reinício, a janela anti-repetição está vazia, mas só são aceites chaves de uma
época posterior: as mensagens da época anterior, incluindo repetições capturadas
antes do reinício, são rejeitadas como repetição, e o tráfego retoma quando os
emissores passam à época seguinte. Um teste de aceitação verifica esta sequência
(\texttt{AfterARestartOnlyALaterEpochIsAccepted}), e outro que uma falha a ler ou
a escrever a época recusa as mensagens em vez de as aceitar. No laboratório, o
\emph{gateway} corre sempre com esta época em disco.
O tempo sem serviço depois de um reinício é, por isso, o tempo que a implantação
leva a passar os emissores à época seguinte. Depende da distribuição de chaves,
fora do âmbito deste trabalho (Secção~\ref{sec:ameacas}), e não foi medido.

\subsection{Rotação e revogação}

Os dois cenários passaram no Raspberry~Pi. Em \texttt{key-rotation}, todas as
chaves do conteúdo são substituídas com o tráfego a correr, sem nenhuma rejeição e
com perda de 0,1\,\%. Em \texttt{revocation}, o sensor~1 é revogado (chaves e
palavra-passe): fica de fora e os outros 49 continuam, com o débito entregue a
descer de 1\,000 para 980~mensagens/s, exatamente a parte do sensor revogado.

\section{Ameaças à validade e limitações}
\label{sec:ameacas}

\subsection{Validade de construção}

\begin{itemize}[leftmargin=*]
  \item \textbf{$S(c)$ conta ataques, não a sua gravidade.} Com os pesos usados,
    todos os ataques valem o mesmo; a exclusão por caso de uso compensa em parte,
    porque uma propriedade obrigatória em falta não é trocada por custo.
  \item \textbf{A escuta no \emph{broker} mede um subscritor estranho.} Na matriz
    de ataques, quem escuta é um cliente sem credenciais, que K1 já impede. Um
    \emph{broker} curioso ou comprometido lê o conteúdo em K1, K2 e K9, que não
    cifram perante o \emph{broker}, e $S(c)$ não o mostra. Quem o mostra é o
    filtro por caso de uso (a telemetria sensível exige o modo de cifra) e o
    cenário que compara o que um subscritor lê com TLS e com o envelope.
  \item \textbf{Um ataque travado no \emph{broker} conta como impedido.} As
    configurações com transporte protegido impedem muitos ataques por o atacante
    não chegar ao \emph{broker}; isso deixa de valer se o \emph{broker} for o
    adversário. K10 e os cenários com \emph{broker} ativo medem o que resta sem
    essa barreira.
  \item \textbf{\emph{Broker} ativo só em MQTT sem TLS.} O intermediário não
    termina TLS; um \emph{broker} malicioso real fá-lo-ia. O que o envelope deteta
    não depende do transporte, pelo que o resultado se mantém, mas não foi medido
    com TLS.
\end{itemize}

\subsection{Validade interna}

\begin{itemize}[leftmargin=*]
  \item \textbf{O instrumento pesa na medição.} O tempo de processador do
    \emph{gateway} inclui a contabilidade das medições, e a sua memória inclui as
    amostras de latência da janela e a tabela de mensagens vistas, que cresce nos
    primeiros minutos de cada laboratório. As diferenças entre configurações são
    válidas; os valores absolutos de memória são um majorante.
  \item \textbf{Tempo de ligação.} Varia em degraus de cerca de 100~ms por esperas
    internas da biblioteca cliente; usa-se o tempo de processador por ligação.
  \item \textbf{Relógios.} A latência entre máquinas depende da correção do
    relógio dos sensores, com incerteza igual a metade do melhor tempo de ida e
    volta observado, no máximo 0,21~ms na campanha (0,18~ms nas execuções de
    custo). É da ordem das diferenças de latência mediana entre configurações
    vizinhas (K0 122~$\mu$s, K8 179~$\mu$s), mas a mesma para todas: como a ordem é
    aleatória, desloca as medidas sem favorecer nenhuma configuração. A deriva é
    corrigida (Secção~\ref{sec:metodo}); uma alteração do ritmo do relógio a meio
    de uma execução não o é, e por isso o NTP esteve desligado durante a campanha.
    O desvio entre as duas máquinas cresceu de forma regular até 289~ms (cerca
    de 16~ppm), sempre corrigido.
  \item \textbf{Temperatura e frequência.} Nos três pilotos não houve
    \emph{throttling} (\texttt{throttled} sempre \texttt{0x0}), a frequência
    esteve sempre em 2\,400~MHz e a temperatura entre 44 e 52~°C. O mesmo na
    campanha: \texttt{0x0} e 2\,400~MHz em todas as execuções, temperatura entre
    44,6 e 52,4~°C (o máximo na subida de carga). A ordem aleatória distribui o
    efeito que houver pelas configurações.
  \item \textbf{Sensores simulados.} Os sensores são processos Python num
    computador, não dispositivos restritos: o custo no emissor não foi medido no
    nível~2 (o nível~1 mede envio e receção).
  \item \textbf{Uma só máquina de carga.} Os dois agentes de sensores e o atacante
    correm no mesmo computador. Ele gerou 80\,000~mensagens/s nos pilotos sem
    ficar aquém da taxa pedida, mas o que acontece nele (escalonamento, recolha de
    memória do Python) entra na latência medida.
  \item \textbf{Picos de latência.} Em todas as configurações, incluindo K0, que
    não tem criptografia, surgem nos pilotos picos de 22 a 40~ms no p99,9 e no
    máximo, a cargas intermédias. Por existirem em K0, vêm do ambiente e não das
    camadas medidas; não afetam o p99 abaixo de cerca de 300~mensagens/s por
    sensor. Na campanha, os maiores valores nas execuções de custo foram de
    0,5~ms em K0. A partir de 400~mensagens/s por sensor
    (20\,000~mensagens/s), a latência mediana sobe de cerca de 0,2 para 7--14~ms
    em todas as configurações, sem perda, com o \emph{gateway} entre 40\,\% (K0)
    e 63\,\% (K8) de processador, longe da saturação, e volta a descer quando a
    carga aumenta. O salto acontece à mesma taxa por sensor em todas as
    configurações; foi registado, mas a sua origem não foi determinada. Afeta a
    latência acima de
    $R_{\text{alvo}}$ (200~mensagens/s por sensor) e fica abaixo de $L_{\max}$,
    pelo que não muda nenhum veredicto, mas a latência na subida de carga acima
    desse ponto é um majorante.
  \item \textbf{Serviços alheios.} O Docker esteve ativo no Raspberry~Pi durante o
    teste de fumo e os pilotos, e foi parado para a campanha, tal como o ambiente
    gráfico; os pilotos só serviram para fixar parâmetros, não entram nos
    resultados.
  \item \textbf{O próprio laboratório tinha defeitos.} Três foram encontrados nas
    duas primeiras campanhas completas e corrigidos antes da terceira, a única
    usada (Secção~\ref{sec:desvios}). O mais grave fazia contar como impedidos
    ataques que não chegaram a ser feitos; foi detetado por o ataque de inundação
    aparecer impedido em K7, que não tem limitação. Na campanha usada,
    verificou-se em todas as células dos ataques internos que o atacante se ligou
    como o dispositivo comprometido e enviou mensagens.
\end{itemize}

\subsection{Validade externa}

\begin{itemize}[leftmargin=*]
  \item Um só modelo de \emph{hardware} restrito (Raspberry~Pi~5), um só
    \emph{broker} (Mosquitto) e uma rede local. Um microcontrolador sem sistema
    operativo, ou uma rede com perdas, podem ordenar as configurações de outra
    forma.
  \item \textbf{Sem AES em \emph{hardware}} foi emulado desligando extensões no
    OpenSSL do mesmo processador, o que não reproduz a memória nem a
    \emph{cache} de um processador mais simples.
  \item Carga sintética: no nível~2, mensagens de 16~octetos de conteúdo a taxa
    constante, QoS~0, 50 sensores; o tamanho só foi variado no nível~1. Com
    mensagens maiores, o peso do envelope e a vantagem do AES em \emph{hardware}
    crescem (Secção~\ref{sec:aceleracao}).
\end{itemize}

\subsection{Limitações do protótipo}

As limitações da \PUm{Secção~6.11} mantêm-se, com uma alteração: o reinício do
recetor deixou de reabrir a janela anti-repetição (Secção~\ref{sec:prototipo}), à
custa de exigir uma rotação de chaves depois de cada reinício. Continuam por
resolver a distribuição de chaves fora do laboratório, o armazenamento de chaves
(só em memória), a supressão e o atraso de mensagens por um \emph{broker} ativo, e
a confidencialidade e a prova de origem em simultâneo (cifra seguida de
assinatura).

\section{Conclusões}
\label{sec:conclusoes}

\PorFazer{Responder a cada pergunta numa ou duas frases, com números:
QP1 (custo de cada camada), QP2 (o que cada camada impede), QP3 (tamanho e
saturação), QP4 (aceleração por \emph{hardware}), QP5 (configuração recomendada por
caso de uso).}

\PorFazer{O que o trabalho acrescenta à lacuna L5 da Parte~1 (falta de medições do
custo e do benefício de cada camada, isoladamente e em combinação, em
\emph{hardware} restrito).}

\PorFazer{Trabalho futuro: cifra seguida de assinatura; distribuição de chaves;
medição no emissor em \emph{hardware} restrito; \emph{broker} ativo com TLS;
energia.}


\clearpage
\phantomsection
\addcontentsline{toc}{section}{\refname}
\bibliographystyle{../IEEEtran-pt}
\bibliography{../referencias}

\clearpage
\appendix
\renewcommand{\thesection}{\Alph{section}}
\section{Reprodução dos resultados}
\label{sec:reproducao}

O repositório contém o código, os dados brutos de cada campanha
(\texttt{results/}) e os programas de análise. Os passos para repetir as medições
num Raspberry~Pi estão em \texttt{docs/raspberry-pi.md}; o procedimento e o
significado de cada ficheiro, em \texttt{docs/evaluation.md}.

\begin{itemize}[leftmargin=*]
  \item \textbf{Refazer as tabelas e figuras a partir dos dados brutos:}\\
    \texttt{tools/analysis/analyze.py results/<campanha>} e, nesta pasta,
    \texttt{make dados PI=<campanha>} seguido de \texttt{make}.
  \item \textbf{Repetir as medições:} \texttt{loadgen.py campaign} com a mesma
    semente e o mesmo \texttt{params.toml} (guardados em cada campanha).
\end{itemize}

\PorFazer{Tabela com as campanhas usadas: nome, plataforma, data, \emph{commit},
semente, número de execuções.}

\subsection*{Plataforma de referência}

\geradoRef{desktop/summary}
\geradoRef{desktop/criterion}


\end{document}
